\section{Konklusjon}
Laboratorieøvelsen har gitt god innsikt i hvordan dioder fungerer i praktiske kretser, spesielt som likerettere og i klippekretser. Det ble observert at dioden effektivt kan brukes til å omdanne vekselspenning til likespenning ved å slippe gjennom strøm i én retning, noe som danner grunnlaget for både énveis og toveis likeretting.
\\[1em]
Videre viste forsøket hvordan dioder og zenerdioder oppfører seg i klippekretser. Vanlige dioder begrenser spenningen i fremoverretning til omtrent terskelspenningen, mens zenerdioder også kan begrense spenningen i reversretning ved en gitt zenerspenning. Dette gjør zenerdioder spesielt egnet i applikasjoner der det er behov for spenningsregulering eller beskyttelse mot overspenning.
\\[1em]
Det ble også tydelig at valg av diode avhenger av bruksområde: vanlige dioder benyttes hovedsakelig til likeretterformål, mens zenerdioder brukes når det er behov for å stabilisere eller begrense spenning i begge retninger.
\\[1em]
Bruk av kondensator i likeretterkretser viste seg å være effektivt for å redusere rippel og gi en mer stabil utgangsspenning. Økt kapasitans førte til bedre glatting, noe som stemmer godt med teoretiske sammenhenger.
\\[1em]
Sammenligning av énveis og toveis likeretting viste at toveis likeretter gir høyere gjennomsnittlig utgangsspenning og mindre rippel, ettersom begge halvperioder av inngangssignalet utnyttes. Ulempen er økt kompleksitet og ekstra spenningsfall som følge av flere dioder i ledende tilstand.
\\[1em]
Resultatene samsvarer i stor grad med teorien, men mindre avvik ble observert. Disse kan forklares med ikke-ideell komponentoppførsel, måleusikkerhet og påvirkning fra det praktiske oppsettet. Den ideelle diodemodellen gir en god tilnærming i enkle analyser, men for mer nøyaktige beregninger må man ta hensyn til faktorer som fremspenning og interne tap.
\\[1em]
Samlet sett har forsøket bidratt til økt forståelse av både teoretiske prinsipper og praktiske begrensninger ved bruk av dioder i elektroniske kretser.